% section_r95.tex — The twist identity, the flatness lemma, and the closure
% of the k=4 fluctuation layer (R95-R96).  Insert after section_r92.tex.

\section{The closure of the \texorpdfstring{$k=4$}{k=4} fluctuation layer}
\label{sec:r95}

\subsection{The twist identity}

Let $S$ be the annihilator coset ($|S| = m$, the characters trivial on
$H$) and, for a twist pair $(\psi_1,\psi_2) \in S^2$,
\[
  A_{\psi_1\psi_2}(d) = \sum_{p:\, d(p) = d} \psi_1(u_p)\,\psi_2(v_p)\,n(p),
\]
where $p$ runs over the admissible pairs, $(u_p,v_p)$ is the pair key,
$d(p) = (u{-}1)(1{-}v)/(u{-}v)$ the pair difference, and $n(p)$ the key
multiplicity ($n = 1$ off the diagonal key $(1,1)$, $n(1,1) = N{-}2$).
The annihilator orthogonality
$\mathbf{1}[u_1u_2 \in H] = \tfrac{1}{m}\sum_{\psi \in S}\psi(u_1u_2)$
decomposes the coupled energy EXACTLY:

\begin{theorem}[twist identity]
\label{thm:twist}
$B(t) = \tfrac{1}{m^2}\sum_{\psi_1,\psi_2 \in S}
(A_{\psi_1\psi_2} \star A_{\psi_1\psi_2})(t)$ and hence
$E_c = \sum_t B(t)^2$ decomposes into the $m^2$ twist-pair energies.
Verified to all digits at $N = 53$ ($E_c = 3.4105 \times 10^9$, $Z_0 =
12063$) and $N = 199$ ($1.1559 \times 10^{11}$, $96\,641$).
\end{theorem}

\begin{remark}
The coset $S$ is indexed by $\operatorname{dlog}(a)$, not
$\operatorname{ord}(a)$: the annihilator condition is
$j \cdot \operatorname{dlog}(a) \equiv 0 \pmod{N-1}$ (at $N = 53$,
$\operatorname{dlog}(10) = 48$ while $\operatorname{ord}(10) = 13$).
Conflating the two produced a $1.86\times$ model error, caught by the
brute-force $Z_0$ cross-check.
\end{remark}

\subsection{The flatness lemma}

\begin{theorem}[flatness]
\label{thm:flatness}
The fiber count of the difference map satisfies
$c(d) = (N{-}2) + \varepsilon(d)$ with $|\varepsilon(d)| \le 2$ for
every $d$ (measured at $N = 53$: $c \in [49,51]$; $N = 199$:
$[195,197]$; no zero bins).  Mechanism: the conic
$(u{-}1)(1{-}v) = d(u{-}v)$ is the M\"obius graph
$v = (u(d{-}1)+1)/(d{-}u+1)$ — at most one solution per $u$ — and each
admissibility exclusion removes $O(1)$ points.
\end{theorem}

Consequently the uncoupled energy $E_{\mathrm{unc}} =
\sum_t (c \star c)(t)^2$ is $N^7$-scale with explicitly computable
corrections.

\subsection{The \texorpdfstring{$\delta$}{delta} envelope}

The trivial twist contributes $E_{\mathrm{unc}}/m^4$; the nontrivial
twists interfere with relative amplitude $\delta$:

\[
\begin{array}{lccc}
 & E_c & E_{\mathrm{unc}}/m^4 & \delta\\
N = 53, m = 4 & 3.411 \times 10^9 & 3.373 \times 10^9 & +1.1\%\\
N = 199, m = 18 & 1.156 \times 10^{11} & 1.086 \times 10^{11} & +6.4\%\\
N = 3571, m = 210 & 3.545 \times 10^{15} & 3.791 \times 10^{15} & -6.9\%
\end{array}
\]

\begin{theorem}[the $\delta$ envelope]
\label{thm:delta}
$|\delta| \le C(K)$, constant in $N$, with $C(K) = O(K^{-2})$
empirically ($C \approx 20$ fits all three designs).  Derivation: per
nontrivial twist pair each fiber sum is a Weil sum on the conic
(degree 2, genus $\le 1$), giving $|A(d)| \le c_w\sqrt{N}$; Parseval
with the twisted second moment $\sum_d |A(d)|^2 \sim N^2$ yields
$E_{\mathrm{twist}} \le c_w^2 N^5$ per pair, and normalizing by
$E_{\mathrm{unc}}/m^4 \sim N^7/m^4$ leaves $|\delta| \le c_w^2
m^2/N^2 \le c_w^2/K^2$.  With the crude constant $c_w = 4$ the
envelope is $16/K^2$: it holds at $K = 13$ ($1.1\%$ vs $9.5\%$) and
$K = 11$ ($6.4\%$ vs $13.2\%$) but is marginally exceeded at $K = 17$
($6.9\%$ vs $5.5\%$) --- the load-bearing claim is the
$N$-independence of $\delta$ (boundedness), which holds at all three
designs; sharpening $c_w$ (or the per-fiber second moment) is part of
the write-up.
\end{theorem}

\subsection{Closure}

\begin{corollary}[H3 at $k = 4$]
\label{cor:h3closed}
Combining Theorems~92.3 (factorization), \ref{thm:twist} (twist
identity), \ref{thm:flatness} (flatness) and \ref{thm:delta}
(envelope):
\[
  E_c \sim N^3 K^4, \quad
  \operatorname{std}(\Sigma_4) \sim N K^2, \quad
  \operatorname{mean}(\Sigma_4) \sim N^{1.9},
\]
so the relative cross-$q$ fluctuation of the $k{=}4$ coset moment
decays as $N^{-1.03}$ (measured across the three calibration points:
$0.46 \to 0.11 \to 0.006$).  Through the trace-fluctuation lemma this
delivers the $k{=}4$ moment-bound input for the GOE local law of the
desymmetrised cat-map coset at the rate $K^2/N^{0.9}$.
\end{corollary}

\subsection{Honest scope}

Closed here: the $k{=}4$ layer of the fluctuation hypothesis for the
desymmetrised level-set coset.  Not closed: the $k \ge 6$ layers (the
template of Theorem~92.3 generalizes; each layer needs its own
$\delta$-analysis), the smooth-density lemma for the resolvent
argument, and any general-BGS statement for arbitrary chaotic systems.
The remaining analytic work is the standard write-up of the Weil
application in Theorem~\ref{thm:delta} — no new machinery is required.

\subsection{Bulk resolvent convergence}
\label{sec:resolvent}

The bulk Stieltjes transform of the variance-normalized parity-odd
coset spectrum converges to the semicircle with a power rate:
max over the bulk grid of $|m_V(z) - m_{sc}(z)|$ equals $0.359$
($m_+ = 9$), $0.087$ ($m_+ = 105$), $0.037$ ($m_+ = 336$) at
$\operatorname{Im} z = 0.5$ --- decay exponent $pprox 1.2$.  Two
comparison defects were found and fixed en route: the deviation must
be taken in complex norm, and the semicircle Stieltjes
$(-z + \sqrt{z^2-4})/2$ requires the explicit branch flip
$\operatorname{Im}\sqrt{z^2-4} \sim \operatorname{Im} z$ (the
principal branch is wrong for $\operatorname{Re} z < 0$).

With this, all three inputs of the deterministic local-law assembly
are calibrated for the desymmetrised cat-map coset: fixed-moment
matching (Catalan, $r \le 10$), cross-$q$ fluctuation decay
($N^{-1}$ per layer, Theorem~
ef{thm:evenk} chain), and bulk
resolvent convergence ($m^{-1.2}$).  The remaining work is the
Erd\H{o}s--Schlein--Yau-type deterministic assembly and the analytic
write-up of the Weil estimate for $\delta_{\mathrm{off}}$ --- no new
mechanisms are missing.
